\QHeading{第1题　口算。}
\par\noindent


\begin{oralgrid}
\OralItem{1}{200-93}{107}
\OralItem{2}{800-350}{450}
\OralItem{3}{150-110}{40}
\OralItem{4}{105-80}{25}
\OralItem{5}{99-27}{72}
\OralItem{6}{280-56}{224}
\OralItem{7}{597-97}{500}
\OralItem{8}{64-39}{25}\end{oralgrid}
\par

\QHeading{口算思路补充（对应上表题号）}

\Answer{上面列出得数，下面按小题给出不同计算思路。计数单位法适用于小数，整数题可比较凑整、拆分或商不变等方法。}
\MethodTitle{第1小题：$200-93$}

\MethodBlock{口算思路一：补差法}{从减数先补到一个整数，再补到被减数，两段差相加。}{\begin{aligned}[t]&\phantom{={}}200-93\\
&=(100-93)+(200-100)\\
&=7+100\\
&=107\end{aligned}}
\MethodTitle{第2小题：$800-350$}

\MethodBlock{口算思路一：补差法}{从减数先补到一个整数，再补到被减数，两段差相加。}{\begin{aligned}[t]&\phantom{={}}800-350\\
&=(360-350)+(800-360)\\
&=10+440\\
&=450\end{aligned}}
\MethodTitle{第3小题：$150-110$}

\MethodBlock{口算思路一：补差法}{从减数先补到一个整数，再补到被减数，两段差相加。}{\begin{aligned}[t]&\phantom{={}}150-110\\
&=(120-110)+(150-120)\\
&=10+30\\
&=40\end{aligned}}
\MethodTitle{第4小题：$105-80$}

\MethodBlock{口算思路一：补差法}{从减数先补到一个整数，再补到被减数，两段差相加。}{\begin{aligned}[t]&\phantom{={}}105-80\\
&=(90-80)+(105-90)\\
&=10+15\\
&=25\end{aligned}}
\MethodTitle{第5小题：$99-27$}

\MethodBlock{口算思路一：补差法}{从减数先补到一个整数，再补到被减数，两段差相加。}{\begin{aligned}[t]&\phantom{={}}99-27\\
&=(30-27)+(99-30)\\
&=3+69\\
&=72\end{aligned}}
\MethodTitle{第6小题：$280-56$}

\MethodBlock{口算思路一：补差法}{从减数先补到一个整数，再补到被减数，两段差相加。}{\begin{aligned}[t]&\phantom{={}}280-56\\
&=(60-56)+(280-60)\\
&=4+220\\
&=224\end{aligned}}
\MethodTitle{第7小题：$597-97$}

\MethodBlock{口算思路一：补差法}{从减数先补到一个整数，再补到被减数，两段差相加。}{\begin{aligned}[t]&\phantom{={}}597-97\\
&=(100-97)+(597-100)\\
&=3+497\\
&=500\end{aligned}}

\MethodBlock{口算思路二：多减再加回}{}{\begin{aligned}[t]&\phantom{={}}597-97\\
&=597-100+3\\
&=497+3\\
&=500\end{aligned}}
\MethodTitle{第8小题：$64-39$}

\MethodBlock{口算思路一：补差法}{从减数先补到一个整数，再补到被减数，两段差相加。}{\begin{aligned}[t]&\phantom{={}}64-39\\
&=(40-39)+(64-40)\\
&=1+24\\
&=25\end{aligned}}

\MethodBlock{口算思路二：多减再加回}{}{\begin{aligned}[t]&\phantom{={}}64-39\\
&=64-40+1\\
&=24+1\\
&=25\end{aligned}}


\par\noindent

\QHeading{第2题　在\(\bigcirc\)和\(\square\)里分别填写相应的运算符号和数。}

\PlainItem{\(2100-997=2100-(\){\color{answerred}1000}\(-\){\color{answerred}3}\()\)\(359-83-59=359-\)}
\par

\PlainItem{{\color{answerred}59}\(-83\)}
\PlainItem{\(264-(64+50)=264-64-\){\color{answerred}50}\(4520+978-520=4520-520+\)}
\PlainItem{{\color{answerred}978}}
\PlainItem{\(526-27-73=526-(\){\color{answerred}27}\(+\){\color{answerred}73}\()\)\(156-68-32=\)}
\PlainItem{{\color{answerred}156}\(-(68\){\color{answerred}\(+\)}\(32)\)}
\par\noindent

\QHeading{第3题　用简便方法计算下面各题。}

\CalcBlock[方法一：简便计算]{1}{\begin{aligned}[t]&\phantom{={}}\color{black}532-223-77\\&\color{answerred}=532-(223+77)\\&\color{answerred}=532-300\\&\color{answerred}=232\end{aligned}}
\par
\MethodBlock{方法二：按运算顺序直接计算（对照）}{本法只作结果核对；本题要求简便计算，前面的简便解法仍保留。}{\begin{aligned}[t]&\phantom{={}}532-223-77\\
&=309-77\\
&=232\end{aligned}}


\CalcBlock[方法一：简便计算]{2}{\begin{aligned}[t]&\phantom{={}}\color{black}485-(185+50)\\&\color{answerred}=485-185-50\\&\color{answerred}=300-50\\&\color{answerred}=250\end{aligned}}
\MethodBlock{方法二：按运算顺序直接计算（对照）}{本法只作结果核对；本题要求简便计算，前面的简便解法仍保留。}{\begin{aligned}[t]&\phantom{={}}485-(185+50)\\
&=485-235\\
&=250\end{aligned}}

\CalcBlock[方法一：简便计算]{3}{\begin{aligned}[t]&\phantom{={}}\color{black}857-287-213\\&\color{answerred}=857-(287+213)\\&\color{answerred}=857-500\\&\color{answerred}=357\end{aligned}}
\MethodBlock{方法二：按运算顺序直接计算（对照）}{本法只作结果核对；本题要求简便计算，前面的简便解法仍保留。}{\begin{aligned}[t]&\phantom{={}}857-287-213\\
&=570-213\\
&=357\end{aligned}}

\CalcBlock[方法一：简便计算]{4}{\begin{aligned}[t]&\phantom{={}}\color{black}767-158-367\\&\color{answerred}=767-367-158\\&\color{answerred}=400-158\\&\color{answerred}=242\end{aligned}}
\MethodBlock{方法二：先合并要减去的部分}{连续减去几个数，可以先把这些数相加再减。}{\begin{aligned}[t]&\phantom{={}}767-158-367\\
&=767-(158+367)\\
&=767-525\\
&=242\end{aligned}}

\MethodBlock{方法三：按运算顺序直接计算（对照）}{本法只作结果核对；本题要求简便计算，前面的简便解法仍保留。}{\begin{aligned}[t]&\phantom{={}}767-158-367\\
&=609-367\\
&=242\end{aligned}}

\CalcBlock[方法一：简便计算]{5}{\begin{aligned}[t]&\phantom{={}}\color{black}563-144-56-263\\&\color{answerred}=(563-263)-(144+56)\\&\color{answerred}=300-200\\&\color{answerred}=100\end{aligned}}
\MethodBlock{方法二：先合并要减去的部分}{连续减去几个数，可以先把这些数相加再减。}{\begin{aligned}[t]&\phantom{={}}563-144-56-263\\
&=563-(144+56+263)\\
&=563-(200+263)\\
&=563-463\\
&=100\end{aligned}}

\MethodBlock{方法三：交换、结合后先凑整}{保留每个数原来的加减号，先把能凑整的部分放在一起。}{\begin{aligned}[t]&\phantom{={}}563-144-56-263\\
&=(563-263)-144-56\\
&=300-144-56\\
&=156-56\\
&=100\end{aligned}}

\MethodBlock{方法四：按运算顺序直接计算（对照）}{本法只作结果核对；本题要求简便计算，前面的简便解法仍保留。}{\begin{aligned}[t]&\phantom{={}}563-144-56-263\\
&=419-56-263\\
&=363-263\\
&=100\end{aligned}}

\CalcBlock[方法一：简便计算]{6}{\begin{aligned}[t]&\phantom{={}}\color{black}536-99\\&\color{answerred}=536-(100-1)\\&\color{answerred}=536-100+1\\&\color{answerred}=436+1\\&\color{answerred}=437\end{aligned}}
\MethodBlock{方法二：补差法}{从减数先补到一个整数，再补到被减数，两段差相加。}{\begin{aligned}[t]&\phantom{={}}536-99\\
&=(100-99)+(536-100)\\
&=1+436\\
&=437\end{aligned}}

\par\noindent

\QHeading{第4题　解决问题。}

\Problem{市场原有536千克蔬菜，上午卖出82千克，下午卖出118千克，晚上该市场还剩下多少千克蔬菜？}
\par


\MethodTitle{方法一}
\Answer{\Formula{\begin{aligned}[t]&\phantom{={}}536-82-118\\&=536-(82+118)\\&=536-200\\&=336\text{（千克）}\end{aligned}}\par 答：晚上这个市场还剩336千克蔬菜。}

\MethodTitle{方法二：先求上午卖完后的剩余}
\Answer{按卖出的先后顺序连续减。}
\Answer{\Formula{\begin{aligned}[t]&\phantom{={}}536-82-118\\
&=454-118\\
&=336\text{（千克）}\end{aligned}}}
\Answer{答：晚上还剩336千克蔬菜。}
